\cotefile{15}
\scanpage{207}
\uline{Faisons varier} $x$ pt complexe (ou \unsure{pl\textsuperscript{ce}|place|plus} archimédienne). Je pense que \uline{les $\unsure{\xi|\zeta}_{x,\mathbb Z}$ et $=$ les $\xi_{x,\mathbb R}$} \inserted{$[\simeq \xi_{x,\mathbb C}$ muni de $\tau]$} \uline{qui correspondent à des \unsure{pl\textsuperscript{ces}|places} distinctes sont non isomorphes}
% slanted
\margin{\tikz[baseline=(a.base)]\node[draw,ellipse,inner sep=1pt,align=center](a){faux, cf\\groupoïde\\fondamental\\sur $Z$};}

Quant aux $\xi_{x,\mathbb C}=\xi_{x,\infty}$, eux $=$s, \deleted{\unsure{comme|isomorphes}} \inserted{\unsure{comme}} $\unsure{K|k}$ est \deleted{nous avons dit, il sont} alg. clos ils sont isomorphes entre eux, mais pas canoniquement ....

\uline{NB} En car $p>0$, je pense que $T_{\alpha}(\mathbb R)=\varnothing$ pour $\alpha$ grand, ce qui explique qu'il n'y ait pas de phénomènes de ``composantes connexes'' à considérer alors .... ``Ramification en $\mathbb R$'' \ldots
% slanted
\margin{vérifier si il ne peut y avoir des pt à val. dans une ext. tot. réelle ???}

Si

\begin{center}\rule{3cm}{0.5pt}\\[-0.65em]\rule{2.6cm}{0.5pt}\end{center}

Pt motivique de $X$ à valeurs dans un corps (ou un anneau) $k$ : Objets de $T(k)$. Langage de la \uline{fibre d'un motif en un point}

les p\textsuperscript{ts} forment une catégorie fibrée sur $(Sch)$ muni de fpqc, \deleted{\illegible{1}} c'est un \uline{torseur}. Pour $k$ donné, les isom. d'un \deleted{objet} $\xi\in T(k)$ sur un autre \deleted{\unsure{forment|jouent}} les
% archive sheet number: 206

